% Applied_Analysis
\documentclass{jsreport}
\usepackage{soupreamble}
\title{応用解析学特論A}
\date{Last Update：\today}

\begin{document}
\maketitle

\abstract{物理現象の多くは偏微分方程式(PDE)で表される。最終的には偏微分方程式の数値計算(差分法- Finite Difference Method)をPCでしたい！(7 月頃に解けるようになりたい)}

\section{常微分方程式の差分解法}
	まず常微分方程式(ODE)の差分解法を述べる。
%\begin{remark}
	常微分方程式や偏微分方程式では\textbf{初期値問題}と\textbf{境界値問題}とで大きく取り扱いが異なる。
%\end{remark}
\subsection{一階微分の差分方程式の構成法}
	関数$u(x)$の一階微分を考える\footnote{与えられた関数は十分滑らかで微分可能}。関数$u(x)$の一階微分は
\begin{equation*}
	\frac{du}{dx}= \lim_{h\to 0} \frac{u(x+ h)- u(x)}{h}
\end{equation*}
	と定義される。$h$が十分小さいとして以下のように\textbf{前進差分}，\textbf{後退差分}，\textbf{中心差分}を考えることができる\footnote{どれを使うかは問題による。}。
\begin{equation*}
	\frac{du}{dx}\approx 
\begin{cases}
	\frac{u(x+ h)- u(x)}{h}：前進差分\\
	\frac{u(x)- u(u- h)}{h}：後退差分\\
	\frac{u(x+ h)- u(x- h)}{2h}：中心差分
\end{cases}
\end{equation*}

\subsection{二階微分の差分式の構成法}
	二階微分も一階微分と同様に
\begin{align*}
	\frac{d^2u}{dx^2}
	&= \lim_{h\to 0} \frac{u'(x+ \frac{h}{2})- u'(x- \frac{h}{2})}{h}\\
	&\approx \frac{u'(x+ \frac{h}{2})- u'(x- \frac{h}{2})}{h}\\
	&\approx \frac{\frac{u(x+ h)- u(x)}{h}- \frac{u(x)- u(x- h)}{h}}{h}\\
	&= \frac{u(x+ h)- 2u(x)+ u(x- h)}{h^2}
\end{align*}

\subsection{Taylor展開を使った差分式の構成法}
\begin{equation} \label{eq:Taylor展開を使った差分式}
	\frac{d^2u}{dx^2}\Big|_{x= x_i}\approx au(x_i- h_1)+ bu(x_i)+ cu(x_i+ h_2)
\end{equation}
	$u(x_i- h_1), cu(x_i+ h_2)$のTaylor展開は以下のように書ける。
\begin{align*}
	u(x_i- h_1)
	&= u(x_i)- h_1u'(x_i)+ \frac{h_1^2}{2}u''(x_i)- \frac{h_1^3}{6}u'''(x_i)+ \cdots\\
	cu(x_i+ h_2)
	&= u(x_i)- h_1u'(x_i)+ \frac{h_1^2}{2}u''(x_i)- \frac{h_1^3}{6}u''(x_i)+ \cdots
\end{align*}
	ここで$u''(x)$の項が消えないようにうまく定数$a, b, c$の値を決める。満たすべき条件は以下のとおり
\begin{equation*}
\begin{cases}
	\frac{1}{2}(h_1^2a+ h_2^2c)= 1\\
	a+ b+ c= 0\\
	-h_1a+ h_2c= 0
\end{cases}
	\Longleftrightarrow
\begin{cases}
	a= \frac{2}{h_1(h_1+ h_2)}\\
	b= \frac{-2}{h_1h_2}\\
	c= \frac{2}{h_2(h_1h_2)}
\end{cases}
\end{equation*}

\begin{itemize}
	\item $n$階微分を差分近似する場合，最低$n+ 1$個の点での$u$の値が必要。
	\item $n+ 2$個以上の点を用いる場合は，Taylor展開においてより高次の$n+ 1$階以上の微分高の係数が0となるようにする。
\end{itemize}
	このようにより多くの点での$u$を用いて高次の項を0とする差分式はより精度が高い。

%\begin{souexample}
	一階微分$\frac{du}{dx}$を3点$x_{i- 1}, x_i, x_{i+ 1}$で近似することを考える。
\begin{equation*}
	\frac{du}{dx}\Big|_{x= x_i}\approx au_{i- 1}+ bu_i+ cu_{i+ 1}
\end{equation*}
	このとき満たすべき条件は
\begin{equation*}
\begin{cases}
	\frac{1}{2}(h_1^2a+ h_2^2c)= 0\\
	a+ b+ c= 0\\
	-h_1a+ h_2c= 1
\end{cases}
	\Longleftrightarrow
\begin{cases}
	a= \frac{-h_2}{h_1(h_1+ h_2)}\\
	b= \frac{h_2- h_1}{h_1h_2}\\
	c= \frac{h_1}{h_2(h_1+ h_2)}
\end{cases}
\end{equation*}
%\end{souexample}

	逆に1階微分に対する前進差分，後退差分，中心差分をTaylor展開してみる。
\begin{itemize}
	\item 前進差分
\begin{align*}
	\frac{du}{dx}
	&\approx \frac{u(x+ h)- {\color{red}u(x)}}{h}\\
	&= \frac{1}{h}\left(\cancel{u(x)}+ hu'(x)+ \frac{h^2}{2}u''(x)+ \cdots- \cancel{{\color{red}u(x)}}\right)\\
	&= u'(x)+ O(h)
\end{align*}

	\item 中心差分
\begin{align*}
	\frac{du}{dx}
	&\approx \frac{u(x+ h)- u(x- h)}{2h}\\
	&= \frac{1}{2h}\left(2h\cdot u'(x)+ \frac{2h^3}{6}u'''(x)+ \cdots\right)\\
	&= u'(x)+ O(h^2)
\end{align*}
\end{itemize}
	精度$p$：元の微分との誤差が$h^p$に比例する。

\section{境界値問題}
	常微分方程式を差分法で解く
\begin{align*}
\begin{cases}
	\frac{du^2}{d^2x}+ u+ x= 0\ (0< x <1)\\
	u(0)= 0, u(1)= 0：境界条件
\end{cases}
	\Longleftrightarrow
	厳密解：u= \frac{\sin{x}}{\sin{1}}- x
\end{align*}
	このような微分方程式を差分法で解くための方法
\begin{tcolorbox}[title= 差分法の計算手順]
\begin{enumerate}
	\item 領域を計算格子で分割する。
\begin{itemize}
	\item 格子分割数(等分割)：$J$，格子点数$J+ 1$
	\item 格子幅：$h= 1/J$
	\item 格子点番号：$j= 0, 1, 2, \ldots, J$
	\item 各格子点の座標：$x_j= j\cdot h$
	\item 各格子点の$u$の近似値：$u_j\approx u(x_j)$
\end{itemize}
	\item 微分方程式を差分方程式に近似する。
\begin{equation*}
	\frac{d^2 u}{dx^2}\approx \frac{u(x+ h)- 2u(x)+ u(x- h)}{h^2}
\end{equation*}
	より差分方程式は
\begin{equation} \label{eq:差分方程式}
\begin{cases}
	\frac{u_{j+ 1}- 2u_j+ u_{j- i}}{h^2}+ u_j+ j\cdot h= 0\ \ (j= 1, 2, \ldots, J- 1)\\
	境界条件：u_0= 0, u_J= 0
\end{cases}
\end{equation}
	\zcref{eq:差分方程式}は式変形すると$u_{j- 1}+ (h^ 2- 2)u_j+ u_{j+ 1}= -j\cdot h^3$. \\
	この式は$j= 1, 2, \ldots, J- 1$の$(J- 1)$個の点で成り立っている。\\
	$\leadsto$ \zcref{eq:差分方程式}は$u_j$に関する$J- 1$元連立一次方程式
\end{enumerate}
\end{tcolorbox}

\section{初期値問題}
	以下のような微分方程式を考える。
\begin{equation}
\begin{cases}
	\frac{du}{dt}= f(t, u)\ (t> 0で時間を表す。)\\
	u(0)= a：初期条件
\end{cases}
\end{equation}

\begin{tcolorbox}[title= 初期値問題の差分方程式の立て方]
\begin{enumerate}
	\item 領域(定義域)を計算格子で分割
\begin{itemize}
	\item 格子幅：$\Delta t$
	\item 格子点番号：$j= 0, 1, 2, \ldots, $
	\item 格子点$n$における時間；$n\Delta t$
	\item 格子点$n$における$u$の近似値：$u^ n\approx u(n\Delta t)$
\end{itemize}
	\item 微分方程式を差分方程式に近似する。
\begin{equation}
	\frac{du}{dt}\approx \frac{u^{n+ 1}- u^n}{\Delta t}\ (前進差分)
\end{equation}
	であるから，差分方程式：
\begin{equation} \label{eq:初期値問題における差分方程式}
\begin{cases}
	u^{n+ 1}= u^n+ \Delta t\cdot f(n\Delta t, u^n),\\
	u^{(0)}= a
\end{cases}
\end{equation}
	が得られる。
\end{enumerate}
\end{tcolorbox}

\section{Euler法}
	時間に関する一階微分を\underline{\textbf{前進差分近似}}して，漸化式の形にして解く方法

	\zcref{eq:初期値問題における差分方程式}を使って初期条件から出発して$\Delta t$刻みで順次近似解が求まる。

\begin{tcolorbox}
	初期値問題\footnote{Euler法を使った解法では}では「連立一次方程式」は解かない\footnote{境界値問題のように}。$\leadsto$ 漸化式を使って順次近似解を計算する。
\end{tcolorbox}

\subsection{具体的な初期値問題}
\begin{equation}
\begin{cases}
	\frac{du}{dt}= -u(t)\\
	u(0)= 1
\end{cases}
	\Longleftrightarrow
	厳密解：u(t)= e^{-t}
\end{equation}

\begin{tcolorbox}
	差分方程式は$\begin{cases}u^{n+ 1}= (1- \Delta t)u^n\ (n= 0, 1, \ldots),\\ u^0= 1\end{cases}$\\
	Euler法で計算すると
\begin{align*}
	u_0&= L\\
	u_1&= (1- \Delta t)u^0= 1- \Delta t\\
	u_2&= (1- \Delta t)u^1= (1- \Delta t)^2\\
	&\vdots\\
	u_n&= (1- \Delta t)^n\xrightarrow{\Delta t\to 0} e^{-t}
\end{align*}
\end{tcolorbox}

\begin{tcolorbox}[title= 微分方程式の解と差分方程式の解]
	(ここで挙げた初期値問題の例では)「差分方程式の解」は差分区間(格子幅)$\Delta t\to 0$の極限で元の「微分方程式の解」に一致する。
	ただし一般的には差分方程式は差分区間$\to 0$の極限で元の微分方程式に一致するが，「差分方程式の\underline{解}」は差分区間$\to 0$の極限で元の「微分方程式の\underline{解}」に一致するとは限らない！！
\end{tcolorbox}

\section{初期値問題のより高精度な計算方法}
\begin{souexample}
\begin{equation}
\begin{cases}
	\frac{du}{dt}= f(t, u)\ (t> 0)\footnote{$z= f(x(t), y(t))$のとき$\frac{\partial z}{\partial t}= \frac{\partial z}{\partial x}\frac{\partial x}{\partial t}+ \frac{\partial z}{\partial y}\frac{\partial y}{\partial t}$}\\
	u(0)= a\ 初期条件
\end{cases}
\end{equation}
\end{souexample}
	$u(t+ \Delta t)$についてTaylor展開
\begin{align*}
	u(t+ \Delta t)
	&= u(t)+ \Delta t\frac{du(t)}{dt}+ \frac{{\Delta t}^2}{2}\frac{d^2u(t)}{dt^2}+ \frac{{\Delta t}^3}{6}\frac{d^3u(t)}{dt^3}\\
	&= u(t)+ \Delta t f(t, u)+ \frac{{\Delta t}^2}{2}\frac{du(t)}{dt}\cancel{+ \frac{{\Delta t}^3}{6}\frac{d^3u(t)}{dt^3}}\\
	&\simeq u(t)+ \Delta tf+ \frac{{\Delta t}^2}{2}\frac{d^2u(t)}{dt^2}\\
\end{align*}
	$\frac{df}{dt}= \frac{\partial f}{\partial t}\frac{dt}{dt}+ \frac{\partial f}{\partial u}\frac{du}{dt}= f_t+ f_u\cdot f$を使って
\begin{align} \label{eq:**}
	u(t+ \Delta t)
	&\simeq u(t)+ \Delta t\cdot f+ \frac{{\Delta t}^2}{2}(f_t+ f_u\cdot f)\\
	&= u(t)+ \Delta t\left(f+ \frac{\Delta t}{2}(f_t+ f_u\cdot f)\right)
\end{align}

\begin{align*}
	pf(t, u)+ qf(t+ r\Delta t, u+ s\Delta tf(t, u))
	&= pf +q(f+ r\Delta tf_t+ s\Delta tf\cdot f_u+ O(\Delta t)^2)\\
	&\simeq (p+ q)f+ q\cdot \Delta t(rf_t+ sff_u)
\end{align*}
	\zcref{eq:**}と比較して$p+ q= 1, q= \frac{1}{2}, r= s= 1$ととれば\zcref{eq:**}は
\begin{equation}
	u(t+ \Delta t)\simeq u(t)+ \frac{\Delta t}{2}\left(f(t, u)+ f(t+ \Delta t, u+ \Delta tf(y, u))\right)
\end{equation}
	と近似することができる。この近似式を使って計算する方法を\textbf{2次精度のRunge– Kutta法}という\footnote{通常Runge– Kutta法を使うときは4次精度}。

\begin{tcolorbox}
\begin{lstlisting}
#include <stdio.h>
#include <math.h>
#define J 20

int main(void)
{
    double a[J], b[J], c[J], d[J], g[J], s[J];
    double u[J+1], v[J+1], h;
    int j;
    FILE *fp;
    
    h= 1.0/J;
    u[0]= 0.0; u[J]= 0.0; /* Boundary condition*/

    for (j= 1; j< J; j++){
        a[j]= c[j]= 1.0;
        b[j]= h*h- 2.0;
        d[j]= -j* h*h*h;
    }

    c[0]= 0.0; c[J- 1]= 0.0;
    g[0]= 1.0; s[0]= 0.0;
    
    for (j= 1; j< J; j++){
        g[j]= b[j]- (a[j]*c[j- 1])/g[j-1];
        s[j]= d[j]- (a[j]*s[j- 1])/g[j-1];
    }

    for (j= J-1; j> 0; j--)
        u[j]= (s[j]- c[j]*u[j+1])/g[j];

    for (j= 0; j<= J; j++)
    v[j]= sin(j*h)/sin(1.0)- j*h;

    fp= fopen("20260518_bvp.csv", "w");
    fprintf(fp, "x, u, v\n");
    for (j= 0; j<= J; j++)
        fprintf(fp, "%f, %f, %f\n", j*h, u[j], v[j]);
    fclose(fp);

    return 0;
}
\end{lstlisting}
\end{tcolorbox}

\subsection{2階線形偏微分方程式の分類}
	$u= u(x, y)$\ 二次元の場合
	$A\frac{d^2u}{dx^2}+ B\frac{d^2u}{dxdy}+ C\frac{du}{dx}+ D\frac{du}{dy}+ Fu+ G= 0$\ ただし$A, B, C, D, E$は定数または$u(x, y)$の関数。
	このとき2階線形偏微分方程式は
\begin{itemize}
	\item $B^2- 4AC< 0$のとき楕円型
	\item $B^2- 4AC= 0$のとき放物型
	\item $B^2- 4AC> 0$のとき双曲型
\end{itemize}
	という。これらの型によって，偏微分方程式の数学的な性質が異なる$\leadsto$ 型によって数値解法も異なる。
\subsubsection{方程式の例}
\begin{itemize}
	\item 楕円型($A= C= 1, B= 0$)
	Laplace方程式：$\frac{\partial^2 u}{\partial x^2}+ \frac{\partial^2 u}{\partial y^2}= 0$\\
	Poisson方程式：$\frac{\partial^2 u}{\partial x^2}+ \frac{\partial^2 u}{\partial y^2}= -f(x, y)$

	\item 放物型($A= a, B= C= 0$)
	拡散方程式：$\frac{\partial u}{\partial t}= a\frac{\partial^2 u}{\partial x^2}$

	\item 双曲型($A= -C^2, B= 0, C= 1$)
	波動方程式：$\frac{\partial^2 u}{\partial x^2}- C^2\frac{\partial^2 u}{\partial x^2}= 0$\\
	$t\to y$と考え，$v= \frac{\partial u}{\partial t}- c\frac{\partial u}{\partial x}$とおくと移流方程式：$\frac{\partial u}{\partial x}- C^2\frac{\partial u}{\partial x}= 0$
\end{itemize}

\begin{itemize}
	\item 楕円型方程式$\to$ 境界値問題
\begin{enumerate_roman}
	\item 領域を分割する($J\times K$)
	$\Delta x= 1/J, \Delta y= 1/K, x_j= j\Delta x, y_k= k\Delta y, u_{j, k}\cong u(x_j, y_k)$
	\item 微分を差分近似する(中心差分法)
	$\frac{u_{{j+ 1}, k}- 2u_{j, k}+ u_{j- 1, k}}{(\Delta x)^2}+ \frac{u_{j, k+1}- 2u_{j, k}+ u_{j, k- 1}}{(\Delta y)^2}= 0$\\
	これを$u_{j, k}$について解くと，
	$u_{j, k}= \frac{1}{2((\Delta x)^2+ (\Delta y)^2)}((\Delta y)^2(u_{j+ 1, k}+ u_{j- 1, k}+ (\Delta x)^2(u_{j, k+ 1}+ u_{j, k- 1}))$
\end{enumerate_roman}
	\item 放物型方程式
	例：一次元熱伝導方程式(拡散方程式)
	長さ1の針金の温度の\underline{時間変化}
\begin{equation}
	\frac{\partial u}{\partial t}= a\frac{\partial^2 u}{\partial x^2}\ (0< x< 1, t> 0, a> 0)
\end{equation}
	初期条件(ini val)：$u(x, 0)= f(x)$\\
	境界条件(bdd val)：$u(0, t)= u(1, t)= 0$
\begin{enumerate_roman}
	\item 領域を分割する。\\
	$\Delta x= 1/J, \Delta t\ (t方向の格子を取る), x_j= j\Delta x, t_n= n\Delta t, u^n_j\cong u(x_j, t_n)$
	\item 微分を差分近似する。
\begin{itemize}
	\item 時間微分：$\frac{\partial u}{\partial t}\cong \frac{u^{n+ 1}- u^n_j}{\Delta t}= a\frac{u^n_{j+ 1}- 2u^n_j+ u^n_{j- 1}}{(\Delta x)^2}$\ (前進差分)
	\item 空間微分：$\frac{\partial^2 u}{\partial x^2}\cong \frac{u^n_{j+ 1}- 2u^n_j+ u^n_{j- 1}}{(\Delta x)^2}$\ (中央差分)
\end{itemize}
	$u^{n+ 1}_j$について整理すると，
	$u^{n+ 1}_j= ru^n_{j+ 1}+ (1- 2r)u^n_j+ ru^n_{j- 1}\ (j= 1, 2, \ldots, J- 1, n= 0, 1, \ldots, r= \frac{a\Delta t}{(\Delta x)^2})$
\end{enumerate_roman}
\end{itemize}

%% 20260615_Applied_Analysis

    $u^{n+ 1}= Su^n$は波数$\xi$の正弦波で表される特殊解をもつ。
    $u^n_j= g^n\exp{(i\xi(j\Delta x))}= g^n\exp{(i\xi x_j)}\ (初期条件：u^0_j= \exp{(i\xi x_j)})$

    $u^{n+ 1}= Su^n$の解をFourier変換したときの波数$\xi$の成分を取り出したもの。各端数成分の線型結合が一般解
    $\because \tau u^n_j= \tau(g^n\exp{(i\xi x_j)})= g^n\exp{(i\xi(x_j+ \Delta x))}= g^n\exp{(i\xi \Delta x)}\exp{(i\xi x_j)}$
    $\because \tau^mu^n_j= g^n\exp{(i\xi(m\Delta x))}\exp{(i\xi x_j)}$
    $u^{n+ i}= Su^n$に代入すると，
    $g^{n+ 1}e\exp{(i\xi x_j)}= \sum_m c_m\exp{(i\xi m\Delta x)}\exp{(i\xi x_j)}\to g= \sum_m c_m\exp{(i\xi m\Delta x)}$

    このように$g$を決めれば$u^n_j= g^n\exp{(i\xi x_j)}$は$u^{n+ 1}= Su^n$の解となることが確認できた。
    $S$によって時間を$Delta t$進めたときに解$u^n_j$がどれだけ拡大(縮小)されているかを表す。

    一般に任意の関数は端数の異なる正弦波の重ね合わせとして表すことができる。$\Delta x= h(\Delta t)$と表したときLaxの安定条件は
\begin{tcolorbox}
    $0\le n\Delta t\le T$を満たす任意の$n\Delta t$に対して，$|g|^n\le C$となるような正定数$u^{n+ 1}= Su^n$の解は$\Delta t\to 0$のとき$\frac{\partial u}{\partial t}= Lu$の解に収束する。 
\end{tcolorbox}
    $\Leftrightarrow$
\begin{tcolorbox}[title= Von\ Neumannの条件]
    $0\le n\Delta t\le T$を満たす任意の$n\Delta t$に対して，$|g|\le 1+ K\Delta t$となるような正定数$K$が存在するとき$u^{n+ 1}= Su^n$の解は$\Delta t\to 0$のとき$\frac{\partial u}{\partial t}= Lu$の解に収束する。
\end{tcolorbox}
    なお$K= 0$ならば$|g|\le 1$となり，誤差は時間とともに減衰するか一定値で置き換えられる。
    なぜなら$t= T$のとき$n\Delta t= T$とすると$\lim_{\Delta t\to 0} (1+ K\Delta t)^n= \lim_{n\to infty} \left(1+ \frac{KT}{n}\right)^n= e^{KT}= C$とおけば$|g|^n\le C$
    $\because g= \sum_m c_m\exp{(i\xi m\Delta x)}$

    Von Neumannの条件をみたすように$\Delta x= h(\Delta x)$を決めておかないと一般には$\Delta t\to 0$のとき差分方程式の解は$\Delta t\to 0$のとき微分方程式は解に収束しない。

%後でノートとる(写真は残してある)

    一次元拡散方程式の例
    一次元拡散方程式
    \begin{equation*}
    \frac{\partial u}{\partial t}= \kappa\frac{\partial^2 u}{\partial x^2}
\end{equation*}
    をEuler陰解法で特殊解の安定化条件を求めよ
%% 20260615_Applied_Analysis

    $u^{n+ 1}= Su^n$は波数$\xi$の正弦波で表される特殊解をもつ。
    $u^n_j= g^n\exp{(i\xi(j\Delta x))}= g^n\exp{(i\xi x_j)}\ (初期条件：u^0_j= \exp{(i\xi x_j)})$

    $u^{n+ 1}= Su^n$の解をFourier変換したときの波数$\xi$の成分を取り出したもの。各端数成分の線型結合が一般解
    $\because \tau u^n_j= \tau(g^n\exp{(i\xi x_j)})= g^n\exp{(i\xi(x_j+ \Delta x))}= g^n\exp{(i\xi \Delta x)}\exp{(i\xi x_j)}$
    $\because \tau^mu^n_j= g^n\exp{(i\xi(m\Delta x))}\exp{(i\xi x_j)}$
    $u^{n+ i}= Su^n$に代入すると，
    $g^{n+ 1}e\exp{(i\xi x_j)}= \sum_m c_m\exp{(i\xi m\Delta x)}\exp{(i\xi x_j)}\to g= \sum_m c_m\exp{(i\xi m\Delta x)}$

    このように$g$を決めれば$u^n_j= g^n\exp{(i\xi x_j)}$は$u^{n+ 1}= Su^n$の解となることが確認できた。
    $S$によって時間を$Delta t$進めたときに解$u^n_j$がどれだけ拡大(縮小)されているかを表す。

    一般に任意の関数は端数の異なる正弦波の重ね合わせとして表すことができる。$\Delta x= h(\Delta t)$と表したときLaxの安定条件は
\begin{tcolorbox}
    $0\le n\Delta t\le T$を満たす任意の$n\Delta t$に対して，$|g|^n\le C$となるような正定数$u^{n+ 1}= Su^n$の解は$\Delta t\to 0$のとき$\frac{\partial u}{\partial t}= Lu$の解に収束する。 
\end{tcolorbox}
    $\Leftrightarrow$
\begin{tcolorbox}[title= Von\ Neumannの条件]
    $0\le n\Delta t\le T$を満たす任意の$n\Delta t$に対して，$|g|\le 1+ K\Delta t$となるような正定数$K$が存在するとき$u^{n+ 1}= Su^n$の解は$\Delta t\to 0$のとき$\frac{\partial u}{\partial t}= Lu$の解に収束する。
\end{tcolorbox}
    なお$K= 0$ならば$|g|\le 1$となり，誤差は時間とともに減衰するか一定値で置き換えられる。
    なぜなら$t= T$のとき$n\Delta t= T$とすると$\lim_{\Delta t\to 0} (1+ K\Delta t)^n= \lim_{n\to infty} \left(1+ \frac{KT}{n}\right)^n= e^{KT}= C$とおけば$|g|^n\le C$
    $\because g= \sum_m c_m\exp{(i\xi m\Delta x)}$

    Von Neumannの条件をみたすように$\Delta x= h(\Delta x)$を決めておかないと一般には$\Delta t\to 0$のとき差分方程式の解は$\Delta t\to 0$のとき微分方程式は解に収束しない。

%後でノートとる(写真は残してある)

    一次元拡散方程式の例
    一次元拡散方程式
    \begin{equation*}
    \frac{\partial u}{\partial t}= \kappa\frac{\partial^2 u}{\partial x^2}
\end{equation*}
    をEuler陰解法で特殊解の安定化条件を求めよ

\end{document}